\documentclass[14pt,a4paper]{extarticle}
% =============================================================
% preamble.tex — оформление отчётов по СТП БГУИР 01-2024
% =============================================================

% Шрифты: под pdflatex — Times через tempora + T2A,
% под xelatex/tectonic — системный Times New Roman (запасной вариант STIX Two Text)
\usepackage{iftex}
\ifPDFTeX
  \usepackage[utf8]{inputenc}
  \usepackage[T2A]{fontenc}
  \usepackage{tempora}
\else
  \usepackage{fontspec}
  \IfFontExistsTF{Times New Roman}
    {\setmainfont{Times New Roman}}
    {\setmainfont{STIX Two Text}}
  \IfFontExistsTF{DejaVu Sans Mono}
    {\setmonofont{DejaVu Sans Mono}[Scale=0.82]}
    {\setmonofont{DejaVuSansMono}[Extension=.ttf, BoldFont=*-Bold,
                                  ItalicFont=*-Oblique, Scale=0.82]}
\fi
\usepackage[russian]{babel}

% Поля: слева 30 мм, справа 15 мм, сверху и снизу 20 мм
\usepackage{geometry}
\geometry{left=30mm, right=15mm, top=20mm, bottom=20mm}

\usepackage{setspace}
\singlespacing

\pretolerance=1000
\tolerance=2000
\emergencystretch=3em

\usepackage{indentfirst}
\setlength{\parindent}{1.25cm}

% Номер страницы — снизу справа
\usepackage{fancyhdr}
\pagestyle{fancy}
\fancyhf{}
\fancyfoot[R]{\thepage}
\renewcommand{\headrulewidth}{0pt}
\fancypagestyle{plain}{%
  \fancyhf{}%
  \fancyfoot[R]{\thepage}%
  \renewcommand{\headrulewidth}{0pt}%
}

\usepackage{amsmath,amssymb}
\usepackage{graphicx}
\usepackage{float}
\usepackage{booktabs}
\usepackage{tabularx}
\usepackage{multirow}
\usepackage{longtable}
\usepackage{array}

\usepackage{caption}
\captionsetup[figure]{labelsep=endash, justification=centering,
                      font={normalsize}, name=Рисунок}
\captionsetup[table]{labelsep=endash, justification=raggedright,
                     singlelinecheck=false, font={normalsize}, name=Таблица}

\usepackage{titlesec}
\titleformat{\section}
  {\newpage\normalfont\fontsize{14}{18}\selectfont\bfseries}
  {\thesection}{0.5em}{\MakeUppercase}
\titlespacing*{\section}{1.25cm}{0pt}{14pt}
\titleformat{\subsection}
  {\normalfont\fontsize{14}{18}\selectfont\bfseries}
  {\thesubsection}{0.5em}{}
\titlespacing*{\subsection}{1.25cm}{14pt}{14pt}
\titleformat{\subsubsection}
  {\normalfont\fontsize{14}{18}\selectfont\bfseries}
  {\thesubsubsection}{0.5em}{}
\titlespacing*{\subsubsection}{1.25cm}{14pt}{14pt}

\newcommand{\sectioncenter}[1]{%
  \newpage
  \phantomsection
  \addcontentsline{toc}{section}{#1}
  \begin{center}\textbf{\MakeUppercase{#1}}\end{center}
  \vspace{14pt}
}

\usepackage{tocloft}
\renewcommand{\cftsecfont}{\normalfont}
\renewcommand{\cftsecpagefont}{\normalfont}
\renewcommand{\cftsubsecfont}{\normalfont}
\renewcommand{\cftsubsecpagefont}{\normalfont}
\renewcommand{\cftsecleader}{\cftdotfill{\cftdotsep}}
\renewcommand{\cftdotsep}{1}
% babel переопределяет заголовки после преамбулы, поэтому правим через \addto
\addto\captionsrussian{\renewcommand{\contentsname}{\hfill\textbf{СОДЕРЖАНИЕ}\hfill}}
\renewcommand{\cftaftertoctitle}{\hfill}

\usepackage{enumitem}
\setlist[itemize]{noitemsep, topsep=0pt, leftmargin=1.25cm, label=--}
\setlist[enumerate]{noitemsep, topsep=0pt, leftmargin=1.25cm}

% Листинги кода и вывода программы
\usepackage{listings}
\usepackage{fancyvrb}
\lstset{
  basicstyle=\ttfamily\footnotesize,
  frame=single,
  breaklines=true,
  showstringspaces=false,
  inputencoding=utf8,
  extendedchars=true,
  columns=fullflexible,
  keepspaces=true,
  captionpos=t
}
\ifPDFTeX
\lstset{
  literate={а}{{\cyra}}1 {б}{{\cyrb}}1 {в}{{\cyrv}}1 {г}{{\cyrg}}1
           {д}{{\cyrd}}1 {е}{{\cyre}}1 {ж}{{\cyrzh}}1 {з}{{\cyrz}}1
           {и}{{\cyri}}1 {й}{{\cyrishrt}}1 {к}{{\cyrk}}1 {л}{{\cyrl}}1
           {м}{{\cyrm}}1 {н}{{\cyrn}}1 {о}{{\cyro}}1 {п}{{\cyrp}}1
           {р}{{\cyrr}}1 {с}{{\cyrs}}1 {т}{{\cyrt}}1 {у}{{\cyru}}1
           {ф}{{\cyrf}}1 {х}{{\cyrh}}1 {ц}{{\cyrc}}1 {ч}{{\cyrch}}1
           {ш}{{\cyrsh}}1 {щ}{{\cyrshch}}1 {ъ}{{\cyrhrdsn}}1
           {ы}{{\cyrery}}1 {ь}{{\cyrsftsn}}1 {э}{{\cyrerev}}1
           {ю}{{\cyryu}}1 {я}{{\cyrya}}1 {ё}{{\cyryo}}1
           {А}{{\CYRA}}1 {Б}{{\CYRB}}1 {В}{{\CYRV}}1 {Г}{{\CYRG}}1
           {Д}{{\CYRD}}1 {Е}{{\CYRE}}1 {Ж}{{\CYRZH}}1 {З}{{\CYRZ}}1
           {И}{{\CYRI}}1 {Й}{{\CYRISHRT}}1 {К}{{\CYRK}}1 {Л}{{\CYRL}}1
           {М}{{\CYRM}}1 {Н}{{\CYRN}}1 {О}{{\CYRO}}1 {П}{{\CYRP}}1
           {Р}{{\CYRR}}1 {С}{{\CYRS}}1 {Т}{{\CYRT}}1 {У}{{\CYRU}}1
           {Ф}{{\CYRF}}1 {Х}{{\CYRH}}1 {Ц}{{\CYRC}}1 {Ч}{{\CYRCH}}1
           {Ш}{{\CYRSH}}1 {Щ}{{\CYRSHCH}}1 {Ъ}{{\CYRHRDSN}}1
           {Ы}{{\CYRERY}}1 {Ь}{{\CYRSFTSN}}1 {Э}{{\CYREREV}}1
           {Ю}{{\CYRYU}}1 {Я}{{\CYRYA}}1 {Ё}{{\CYRYO}}1
           {—}{{---}}1 {–}{{--}}1 {…}{{\ldots}}1
}
\fi

% Подписи листингов — в том же стиле, что таблицы: «Листинг X.Y — Название»
\captionsetup[lstlisting]{labelsep=endash, justification=raggedright,
                          singlelinecheck=false, font={normalsize}}
\renewcommand{\lstlistingname}{Листинг}
\AtBeginDocument{\numberwithin{lstlisting}{section}}

\lstdefinestyle{python}{language=Python, keywordstyle=\bfseries}
\lstdefinestyle{shell}{language=bash, keywordstyle=\mdseries}
\lstdefinestyle{plain}{language={}, basicstyle=\ttfamily\scriptsize}

% Абзац после листинга должен сохранять отступ 1,25 см (СТП 2.1.1)
\makeatletter
\AtBeginDocument{\let\@endpetrue\@endpefalse}
\makeatother

% Подпись к блоку verbatim: \captionof вне плавающего окружения
% глобально обнуляет \parindent, поэтому отступ восстанавливается вручную
\newcommand{\listingcaption}[2]{%
  \captionof{lstlisting}{#1}\label{#2}%
  \setlength{\parindent}{1.25cm}%
}

\usepackage{hyperref}
\hypersetup{colorlinks=true, linkcolor=black, urlcolor=blue, citecolor=black}

\numberwithin{figure}{section}
\numberwithin{table}{section}
\numberwithin{equation}{section}

% Титульный лист отчёта по лабораторной работе
\newcommand{\makelabtitle}[2]{%
  \begin{titlepage}
  \thispagestyle{empty}
  \begin{center}

  МИНИСТЕРСТВО ОБРАЗОВАНИЯ РЕСПУБЛИКИ БЕЛАРУСЬ

  \vspace{0.5em}
  Учреждение образования

  \vspace{0.5em}
  \textbf{<<Белорусский государственный университет информатики и радиоэлектроники>>}

  \vspace{2em}
  Факультет компьютерных систем и сетей

  \vspace{0.5em}
  Кафедра программного обеспечения информационных технологий
  \end{center}

  \vspace{3em}
  \begin{center}
  \textbf{ОТЧЕТ}

  \vspace{0.5em}
  по лабораторной работе №#1

  \vspace{0.5em}
  по дисциплине

  \vspace{0.5em}
  <<Информационные и технические средства защиты авторских прав>>

  \vspace{0.5em}
  <<#2>>
  \end{center}

  \vspace{5em}
  \begin{flushright}
  \begin{minipage}{0.5\textwidth}
  \textbf{Выполнил:}\\
  Лебедевич Артем Владимирович\\
  магистрант группы 556301

  \vspace{0.6em}
  \textbf{Преподаватель:}\\
  Ярмолик Вячеслав Николаевич\\
  доктор технических наук\\
  профессор кафедры ПОИТ
  \end{minipage}
  \end{flushright}

  \vfill
  \begin{center}Минск, 2026~г.\end{center}
  \end{titlepage}
}

% Титульный лист реферата
\newcommand{\makerefattitle}[1]{%
  \begin{titlepage}
  \thispagestyle{empty}
  \begin{center}

  МИНИСТЕРСТВО ОБРАЗОВАНИЯ РЕСПУБЛИКИ БЕЛАРУСЬ

  \vspace{0.5em}
  Учреждение образования

  \vspace{0.5em}
  \textbf{<<Белорусский государственный университет информатики и радиоэлектроники>>}

  \vspace{2em}
  Факультет компьютерных систем и сетей

  \vspace{0.5em}
  Кафедра программного обеспечения информационных технологий
  \end{center}

  \vspace{3em}
  \begin{center}
  \textbf{РЕФЕРАТ}

  \vspace{0.5em}
  по дисциплине

  \vspace{0.5em}
  <<Информационные и технические средства защиты авторских прав>>

  \vspace{1.5em}
  на тему

  \vspace{0.5em}
  \textbf{<<#1>>}
  \end{center}

  \vspace{5em}
  \begin{flushright}
  \begin{minipage}{0.5\textwidth}
  \textbf{Выполнил:}\\
  Лебедевич Артем Владимирович\\
  магистрант группы 556301

  \vspace{0.6em}
  \textbf{Преподаватель:}\\
  Ярмолик Вячеслав Николаевич\\
  доктор технических наук\\
  профессор кафедры ПОИТ
  \end{minipage}
  \end{flushright}

  \vfill
  \begin{center}Минск, 2026~г.\end{center}
  \end{titlepage}
}


\begin{document}

\makelabtitle{3}{Электронная цифровая подпись RSA}

\tableofcontents

\section{Цель работы}

Изучить назначение и устройство функций хеширования и электронной цифровой подписи.
В соответствии с вариантом №1 (нечётный номер в списке группы) реализовать программу
вычисления и проверки электронной цифровой подписи по схеме RSA, используя для
вычисления хеш-образа сообщения функцию~(3.2) методического пособия. Проверить
корректность реализации на примере из пособия и исследовать стойкость схемы.

\section{Теоретические сведения}

\subsection{Функция хеширования}

Функцией хеширования $h$ называется преобразование, переводящее строку $M$
произвольной длины в значение $m = h(M)$ фиксированной длины --- хеш-образ, или
дайджест сообщения. Криптографическая хеш-функция должна:
\begin{itemize}
\item быть чувствительной к любому изменению входной последовательности;
\item применяться к данным любой длины и давать результат фиксированной длины;
\item быть односторонней: по $h(M)$ невозможно найти $M$;
\item не допускать нахождения второго прообраза $M' \ne M$ с $h(M') = h(M)$;
\item не допускать нахождения произвольной пары $M_1 \ne M_2$ с $h(M_1) = h(M_2)$
(коллизии).
\end{itemize}

Хеш-образ строится поблочно: $H_i = f(H_{i-1}, m_i)$, результат обработки последнего
блока становится хеш-образом всего сообщения. В пособии в качестве учебного примера
используется квадратичная функция
\begin{equation}
H_i = (H_{i-1} + m_i)^2 \bmod n,
\end{equation}
где $n = p \cdot q$, $H_0$ --- произвольное начальное значение, $m_i$ --- код $i$-го
символа: $\text{А} = 1, \text{Б} = 2, \dots, \text{Я} = 33$.

Для сообщения $M = \text{<<БГУИР>>} = \{2, 4, 21, 10, 18\}$ при $n = 17 \cdot 19 = 323$
и $H_0 = 100$ вычисления приведены в таблице~\ref{tab:hash}.

\begin{table}[H]
\caption{Хеш-образ сообщения <<БГУИР>>}
\label{tab:hash}
\centering
\begin{tabular}{cccl}
\toprule
$i$ & Символ & $m_i$ & $H_i = (H_{i-1} + m_i)^2 \bmod 323$ \\
\midrule
1 & Б & 2 & $(100 + 2)^2 \bmod 323 = 68$ \\
2 & Г & 4 & $(68 + 4)^2 \bmod 323 = 16$ \\
3 & У & 21 & $(16 + 21)^2 \bmod 323 = 77$ \\
4 & И & 10 & $(77 + 10)^2 \bmod 323 = 140$ \\
5 & Р & 18 & $(140 + 18)^2 \bmod 323 = 93$ \\
\bottomrule
\end{tabular}
\end{table}

Таким образом, $h(\text{БГУИР}) = 93$.

\subsection{Классическая схема цифровой подписи}

Цифровая подпись играет для электронного документа ту же роль, что подпись от руки
для бумажного: подтверждает, что документ составил или заверил владелец подписи, и
позволяет обнаружить изменение документа после подписания. Схема решает три задачи:
подпись невозможно подделать, её принадлежность субъекту можно проверить, а сам
субъект не может от неё отказаться.

Отправитель, владеющий секретным ключом $K_c$:
\begin{enumerate}
\item вычисляет хеш-образ $m = h(M)$;
\item вычисляет подпись $S$ по хеш-образу с помощью секретного ключа;
\item отправляет пару $(M, S)$.
\end{enumerate}

Получатель пары $(M', S)$, знающий открытый ключ $K_o$:
\begin{enumerate}
\item вычисляет хеш-образ полученного сообщения $m' = h(M')$;
\item извлекает из подписи хеш-образ $m$ с помощью открытого ключа;
\item признаёт подпись подлинной, если $m' = m$.
\end{enumerate}

Подписывается не само сообщение, а его хеш-образ: это делает подпись короткой
независимо от длины документа и избавляет от необходимости выполнять медленное
возведение в степень для каждого блока сообщения.

\subsection{Алгоритм цифровой подписи RSA}

Пара ключей вычисляется так же, как для криптосистемы RSA (лабораторная работа №2):
$r = p \cdot q$, $\varphi(r) = (p - 1)(q - 1)$, открытая экспонента $e$ взаимно проста с
$\varphi(r)$, секретная $d$ удовлетворяет условию $e \cdot d \bmod \varphi(r) = 1$.
Открытый ключ $K_o = (e, r)$ передаётся партнёрам для проверки подписей, секретный
$K_c = (d, r)$ хранит автор.

Роли ключей по сравнению с шифрованием меняются местами: подпись вычисляется
секретным ключом, а проверяется открытым:
\begin{equation}
S = m^{d} \bmod r, \qquad m = S^{e} \bmod r.
\end{equation}

Подпись подлинна, если $h(M') = S^e \bmod r$. Корректность следует из того же
тождества $m^{ed} \equiv m \pmod r$, что и для шифрования. Создать подпись, которая
пройдёт проверку открытым ключом, может только владелец $d$, а вычислить $d$ по
$(e, r)$ так же трудно, как разложить $r$ на множители.

\subsection{Пример из методического пособия}

Для $p = 17$, $q = 19$: $r = 323$, $\varphi(r) = 288$. При $e = 43$ расширенный алгоритм
Евклида даёт $d = 67$ (таблица~\ref{tab:euclid}). Открытый ключ $(43, 323)$,
секретный --- $(67, 323)$.

\begin{table}[H]
\caption{Расширенный алгоритм Евклида для $a = 288$, $b = 43$}
\label{tab:euclid}
\centering
\begin{tabular}{ccrrr}
\toprule
Шаг & $q$ & $d$ & $x$ & $y$ \\
\midrule
0 & --- & 288 & 1 & 0 \\
1 & --- & 43 & 0 & 1 \\
2 & 6 & 30 & 1 & $-6$ \\
3 & 1 & 13 & $-1$ & 7 \\
4 & 2 & 4 & 3 & $-20$ \\
5 & 3 & 1 & $-10$ & 67 \\
\bottomrule
\end{tabular}
\end{table}

Подпись сообщения <<БГУИР>>: $S = 93^{67} \bmod 323 = 206$, отправляется пара
$\{\text{БГУИР}, 206\}$. Пусть при передаче сообщение изменено и получатель принял
$\{\text{БРУИР}, 206\}$. Он вычисляет $h(\text{БРУИР}) = 123$, а из подписи
восстанавливает $206^{43} \bmod 323 = 93$. Значения не совпадают, подпись
признаётся недействительной.

В пособии хеш-функция и подпись используют один и тот же модуль $n = r = 323$.
В программе это сделано правилом: хеш вычисляется по модулю $r$, поэтому любой
хеш-образ является допустимым блоком RSA ($m < r$) при любых ключах.

\section{Реализация}

\subsection{Структура программы}

Код третьей лабораторной работы находится в отдельном пакете \texttt{lab3}. Ключи
RSA вычисляются модулем \texttt{lab2/rsa.py}, а быстрое возведение в степень и
расширенный алгоритм Евклида --- общими функциями \texttt{helpers/numtheory.py}
(листинг~\ref{lst:tree}).

\listingcaption{Структура программы}{lst:tree}
\begin{Verbatim}[fontsize=\footnotesize, frame=single, samepage=true]
main.py                CLI: команды sign и verify, demo 3, gui 3
configs/cfg.py         p, q, e, H0 по умолчанию (пример из пособия)
helpers/numtheory.py   fast_exp, euclid_ex (общие с ЛР2)
lab2/rsa.py            вычисление пары ключей RSA
lab3/hashing.py        квадратичная хеш-функция (формула 3.2)
lab3/signature.py      вычисление и проверка подписи
lab3/runner.py         прогон примеров пособия, коллизия, самопроверка
lab3/gui.py            окно: подписание, проверка, пошаговый разбор
\end{Verbatim}

\subsection{Ключевые фрагменты кода}

Хеш-функция (листинг~\ref{lst:hash}) кодирует символы по таблице пособия
$\text{А} = 1, \dots, \text{Я} = 33$ (с буквой Ё). Чтобы подписывать произвольные
документы, таблица продолжена латинскими буквами ($A = 34, \dots, Z = 59$) и цифрами
($0 = 60, \dots, 9 = 69$). Регистр не учитывается, пробелы и знаки препинания не
имеют кода и пропускаются.

\begin{lstlisting}[style=python, caption={Квадратичная хеш-функция}, label={lst:hash}]
RUSSIAN = "АБВГДЕЁЖЗИЙКЛМНОПРСТУФХЦЧШЩЪЫЬЭЮЯ"
ENGLISH = "ABCDEFGHIJKLMNOPQRSTUVWXYZ"
DIGITS = "0123456789"
CODES = {symbol: i for i, symbol in enumerate(RUSSIAN + ENGLISH + DIGITS, start=1)}


def quad_hash(text, n, h0, trace=None):
    h = h0 % n
    for symbol in (ch for ch in text.upper() if ch in CODES):
        m = CODES[symbol]
        h = (h + m) ** 2 % n
        if trace is not None:
            trace.append(HashStep(symbol, m, h))
    return h
\end{lstlisting}

Вычисление и проверка подписи (листинг~\ref{lst:sign}) занимают несколько строк:
вся сложность сосредоточена в функциях \texttt{fast\_exp} и \texttt{make\_keys},
реализованных и проверенных в лабораторной работе №2. Проверка возвращает оба
сравниваемых значения, чтобы программа могла показать, в чём именно расхождение.

\begin{lstlisting}[style=python, caption={Вычисление и проверка подписи}, label={lst:sign}]
class Verdict(NamedTuple):
    computed: int         # m' = h(M') from the received message
    recovered: int        # m = S^e mod r from the signature
    valid: bool


def sign(text, d, r, h0, trace=None):
    return fast_exp(quad_hash(text, r, h0, trace), d, r)


def verify(text, signature, e, r, h0, trace=None):
    if not 0 <= signature < r:
        raise ValueError(f"The signature must be in 0..r-1 = 0..{r - 1}")
    computed = quad_hash(text, r, h0, trace)
    recovered = fast_exp(signature, e, r)
    return Verdict(computed, recovered, computed == recovered)
\end{lstlisting}

\subsection{Интерфейс командной строки}

Команды \texttt{sign} и \texttt{verify} (листинг~\ref{lst:cli}) принимают ключи либо
через $p$, $q$, $e$, либо напрямую: автору для подписи достаточно $(d, r)$, получателю
для проверки --- открытого ключа $(e, r)$. Флаг \texttt{-v} выводит вычисление
хеш-образа по шагам. При недействительной подписи команда \texttt{verify} завершается
с кодом 1, что позволяет использовать её в скриптах.

\begin{lstlisting}[style=shell, caption={Примеры запуска}, label={lst:cli}]
python3 main.py sign -t БГУИР -v               # хеш по шагам, h(M) = 93, S = 206
python3 main.py sign -d 67 -r 323 -t БГУИР     # то же по секретному ключу
python3 main.py verify -e 43 -r 323 -s 206 -t БГУИР    # VALID
python3 main.py verify -e 43 -r 323 -s 206 -t БРУИР    # INVALID: 123 != 93
python3 main.py demo 3                          # все примеры с проверкой
python3 main.py gui 3                           # графическое окно
\end{lstlisting}

\subsection{Графическое приложение}

Окно лабораторной работы №3 (\texttt{python3 main.py gui 3}) построено на базовом
классе окна лабораторной работы №2 и дополнительно содержит поле $H_0$, поле
сообщения, поле подписи и кнопки <<Подписать>> и <<Проверить>>. При подписании
вычисленная подпись подставляется в поле $S$, на панели разбора выводятся ход
алгоритма Евклида, таблица хеширования и формула $S = m^d \bmod r$ с числами. При
проверке результат выводится цветом: зелёным для подлинной подписи и красным для
недействительной, с обоими сравниваемыми значениями. Это позволяет на занятии
изменить одну букву сообщения или одну цифру подписи и сразу увидеть, как изменился
хеш-образ.

\section{Результаты работы программы}

Команда \texttt{python3 main.py demo 3} выполняет пример из пособия, поиск коллизии,
подпись длинного сообщения 256-битным ключом и самопроверку. Сводные результаты
приведены в таблице~\ref{tab:results}, полный вывод --- в листинге~\ref{lst:demo}.

\begin{table}[H]
\caption{Результаты работы программы}
\label{tab:results}
\centering
\begin{tabularx}{\textwidth}{X>{\raggedright\arraybackslash}p{5.4cm}l}
\toprule
Проверка & Результат & Эталон \\
\midrule
Секретная экспонента, $p = 17$, $q = 19$, $e = 43$ & $d = 67$ & совпадает \\
Хеш-образ <<БГУИР>> & 93 & совпадает \\
Подпись <<БГУИР>> & $S = 206$ & совпадает \\
Проверка $\{\text{БГУИР}, 206\}$ & подлинная & --- \\
Проверка $\{\text{БРУИР}, 206\}$ & $123 \ne 93$, отклонена & совпадает \\
Проверка $\{\text{БГУИР}, 207\}$ & $93 \ne 24$, отклонена & --- \\
Коллизия $\{\text{АБЛ}, 206\}$ & принята как подлинная & --- \\
256-битный ключ, изменена одна цифра & отклонена & --- \\
200 случайных сообщений и ключей & все подлинные & --- \\
Изменение одного символа в них & отклонено в 200 из 200 & --- \\
\bottomrule
\end{tabularx}
\end{table}

\listingcaption{Вывод программы}{lst:demo}
\VerbatimInput[fontsize=\scriptsize, frame=single]{listings/demo_output.txt}

\section{Анализ стойкости}

\subsection{Коллизии хеш-функции при малом модуле}

Стойкость подписи не выше стойкости хеш-функции: подпись связана с сообщением только
через хеш-образ, поэтому любое другое сообщение с тем же хеш-образом примет чужую
подпись как свою. При $n = 323$ хеш-функция принимает не более 323 значений, и
коллизию можно найти простым перебором. Программа перебирает трёхбуквенные слова и
уже на 46-й попытке находит слово <<АБЛ>> с $h(\text{АБЛ}) = 93 = h(\text{БГУИР})$.
Пара $\{\text{АБЛ}, 206\}$ проходит проверку как подлинная, хотя владелец ключа
никогда не подписывал это сообщение.

По парадоксу дней рождения для хеш-образа длиной $k$ бит коллизия находится примерно
за $2^{k/2}$ попыток, поэтому в реальных системах используются хеш-функции длиной не
менее 256 бит (SHA-256, SHA-3, белорусский стандарт СТБ 34.101.31 <<belt-hash>>).

\subsection{Слабости учебной хеш-функции}

Квадратичная хеш-функция пособия не является криптографической и при большом модуле.
Каждый её шаг --- возведение в квадрат, а у квадрата по модулю $n = pq$ четыре
корня: значения $x$ и $n - x$ дают одинаковый результат, поэтому два разных
промежуточных значения $H_{i-1} + m_i$, дающие в сумме $n$, приводят к одному
хеш-образу. Зная разложение $n$, можно извлекать квадратные корни и обращать шаги
хеш-функции. Если хеш вычисляется по модулю $r$ ключа подписи, этим знанием обладает
сам автор: он может заранее построить пару сообщений с одинаковым хеш-образом,
подписать одно и затем утверждать, что подписывал другое. Кроме
того, в реализации пробелы и знаки препинания не участвуют в хешировании, поэтому
сообщения <<БГУ ИР>> и <<БГУИР>> имеют одинаковую подпись. Таблица кодов пособия
определена только для букв, и это ограничение сохранено сознательно.

\subsection{Мультипликативность RSA}

Без хеширования подпись RSA допускает подделку: из $S_1 = m_1^d$ и $S_2 = m_2^d$ любой
получает подпись $S_1 S_2 \bmod r = (m_1 m_2)^d \bmod r$ сообщения $m_1 m_2$, не зная
секретного ключа, а для любого $S$ число $S^e \bmod r$ является <<сообщением>>, которому
$S$ подходит как подпись. Хеширование разрушает эту алгебраическую связь, поскольку
противник должен ещё найти осмысленный документ с нужным хеш-образом. Современный
стандарт RSA-PSS дополнительно добавляет к хеш-образу случайную соль.

\subsection{Цифровая подпись в защите авторских прав}

Для защиты авторских прав цифровая подпись решает две задачи: подтверждает авторство
(подпись под произведением может поставить только владелец секретного ключа) и
целостность (любое изменение произведения после подписания обнаруживается).
В сочетании с меткой времени от доверенной службы подписанный хеш-образ произведения
служит доказательством того, что автор обладал им на определённую дату. Подпись,
однако, не мешает копированию: её можно удалить вместе с файлом-приложением,
поэтому для контроля распространения она дополняется встраиваемыми водяными знаками
и средствами управления цифровыми правами.

\sectioncenter{Заключение}

В ходе лабораторной работы изучены требования к функциям хеширования и классическая
схема электронной цифровой подписи. Реализована программа вычисления и проверки
подписи по схеме RSA (вариант №1) с хеш-функцией~(3.2) пособия, с интерфейсом
командной строки и графическим окном, показывающим вычисления по шагам. Корректность
подтверждена совпадением всех значений с примером пособия ($d = 67$, $h(M) = 93$,
$S = 206$, $h(\text{БРУИР}) = 123$) и проверкой на 200 случайных сообщениях и ключах.

Эксперимент с коллизией показал, что стойкость подписи определяется слабейшим звеном
схемы: при 323 возможных значениях хеш-образа подделка находится за 46 попыток, и
стойкость RSA при этом не играет роли. Надёжная подпись требует одновременно модуля
RSA не менее 2048 бит, криптографической хеш-функции длиной не менее 256 бит и
схемы дополнения хеш-образа, например RSA-PSS.

\sectioncenter{Список использованных источников}

\begin{enumerate}
\item Ярмолик, С.~В. Криптография и охрана коммерческой информации: метод. пособие по
выполнению лаб. работ / С.~В.~Ярмолик, В.~Н.~Ярмолик. --- Минск: БГУИР, 2010. --- 32~с.
\item Rivest, R. A method for obtaining digital signatures and public-key cryptosystems /
R.~Rivest, A.~Shamir, L.~Adleman // Communications of the ACM. --- 1978. --- Vol.~21,
No.~2. --- P.~120--126.
\item Шнайер, Б. Прикладная криптография. Протоколы, алгоритмы, исходные тексты на
языке Си / Б.~Шнайер. --- М.: Триумф, 2002. --- 816~с.
\item Столлингс, В. Криптография и защита сетей: принципы и практика /
В.~Столлингс. --- М.: Вильямс, 2001. --- 672~с.
\end{enumerate}

\end{document}
